\documentclass[12pt, a4paper]{article}
\usepackage[utf8]{inputenc}
\usepackage[T1]{fontenc}
\usepackage[french]{babel}
\usepackage[margin=2.5cm]{geometry}
\usepackage{amsmath, amssymb, amsthm, mathtools}
\usepackage{listings}
\usepackage{xcolor}
\usepackage{tcolorbox}
\tcbuselibrary{theorems, skins, breakable}
\usepackage{booktabs}
\usepackage{fancyhdr}
\usepackage{hyperref}
\usepackage{tikz}
\usetikzlibrary{arrows.meta, positioning, matrix}

\definecolor{suva_blue}{RGB}{30, 100, 180}
\definecolor{code_bg}{RGB}{245, 247, 250}
\definecolor{code_kw}{RGB}{0, 80, 200}
\definecolor{code_str}{RGB}{180, 40, 40}
\definecolor{code_cmt}{RGB}{80, 130, 80}
\definecolor{warn_bg}{RGB}{255, 248, 220}
\definecolor{success_bg}{RGB}{230, 255, 235}

\lstdefinestyle{python}{
  language=Python,
  basicstyle=\ttfamily\small,
  keywordstyle=\color{code_kw}\bfseries,
  stringstyle=\color{code_str},
  commentstyle=\color{code_cmt}\itshape,
  backgroundcolor=\color{code_bg},
  frame=single, rulecolor=\color{gray!40},
  numbers=left, numberstyle=\tiny\color{gray},
  breaklines=true, tabsize=4, showstringspaces=false,
}
\lstdefinestyle{output}{
  basicstyle=\ttfamily\small,
  backgroundcolor=\color{gray!10},
  frame=single, rulecolor=\color{gray!50},
  breaklines=true,
}

\newtcolorbox{defbox}[2][]{colback=success_bg,colframe=green!50!black,
  fonttitle=\bfseries,title={Définition : #2},#1}
\newtcolorbox{theorembox}[2][]{colback=blue!5,colframe=suva_blue,
  fonttitle=\bfseries,title={Théorème : #2},#1}
\newtcolorbox{exercice}[2][]{colback=white,colframe=suva_blue!60,
  fonttitle=\bfseries,title={Exercice #2},#1,breakable}
\newtcolorbox{solution}[1][]{colback=gray!5,colframe=gray!40,
  fonttitle=\bfseries\itshape,title={Solution},#1,breakable}
\newtcolorbox{importantbox}[1][]{colback=suva_blue!10,colframe=suva_blue,
  fonttitle=\bfseries,title={Connexion avec le projet},#1}
\newtcolorbox{warningbox}[1][]{colback=warn_bg,colframe=orange!80!black,
  fonttitle=\bfseries,title={Attention},#1}

\newcommand{\Zq}{\mathbb{Z}_q}
\newcommand{\Rq}{R_q}
\newcommand{\NTT}{\mathrm{NTT}}
\newcommand{\INTT}{\mathrm{NTT}^{-1}}

\pagestyle{fancy}
\fancyhf{}
\fancyhead[L]{\color{suva_blue}\textbf{SuVa --- Chapitre 3}}
\fancyhead[R]{\color{gray}NTT}
\fancyfoot[L]{\footnotesize\color{gray}Mohamed Lamine MBENGUE}
\fancyfoot[C]{\thepage}
\fancyfoot[R]{\footnotesize\color{gray}portfolio-alamine.web.app}
\renewcommand{\headrulewidth}{0.4pt}
\renewcommand{\footrulewidth}{0.3pt}

% ══════════════════════════════════════════════════════════════════════
\begin{document}

\begin{center}
  {\color{suva_blue}\rule{\linewidth}{2pt}}\\[0.5cm]
  {\Huge\bfseries\color{suva_blue} Chapitre 3}\\[0.3cm]
  {\large\bfseries Number Theoretic Transform (NTT)}\\[0.2cm]
  {\small\color{gray} Le moteur de performance de Dilithium}\\[0.3cm]
  {\color{suva_blue}\rule{\linewidth}{2pt}}
\end{center}

\begin{tcolorbox}[colback=suva_blue!5, colframe=suva_blue!40,
  left=6pt, right=6pt, top=4pt, bottom=4pt]
\small
\begin{tabular}{ll}
  \textbf{Auteur} & Mohamed Lamine MBENGUE \\
  \textbf{Spécialités} & Sécurité · DevOps · Dev Full Stack Web \& Mobile \\
  \textbf{Contact} & +221 78 178 44 36 · mbenguemohamedlamine2022@gmail.com \\
  \textbf{Portfolio} & \href{https://portfolio-alamine.web.app}{\texttt{portfolio-alamine.web.app}} \\
\end{tabular}
\end{tcolorbox}

\tableofcontents
\newpage

%══════════════════════════════════════════════════════════════════════
\section{Problème : multiplier des polynômes efficacement}
%══════════════════════════════════════════════════════════════════════

Dans Dilithium, on doit calculer $w = A \cdot y$ : une multiplication
matrice-vecteur de polynômes. Avec $k = l = 4$ et $n = 256$, c'est~:
$$4 \times 4 \times 256^2 = 1\,048\,576 \text{ multiplications en méthode naïve}$$

Avec la NTT, on réduit à~:
$$4 \times 4 \times (256 \log_2 256) = 4 \times 4 \times 2048 = 32\,768$$

Gain~: \textbf{facteur 32}.

%══════════════════════════════════════════════════════════════════════
\section{La DFT --- intuition}
%══════════════════════════════════════════════════════════════════════

La DFT (Discrete Fourier Transform) transforme un signal $f$ de longueur $n$~:
$$\hat{f}_k = \sum_{j=0}^{n-1} f_j \cdot \omega^{jk}, \quad k = 0, \ldots, n-1$$
où $\omega = e^{2\pi i / n}$ est la racine primitive $n$-ième de l'unité.

\textbf{Propriété clé~:} La multiplication devient une multiplication terme-à-terme~:
$$\widehat{f \cdot g}_k = \hat{f}_k \cdot \hat{g}_k$$

Donc~: $f \cdot g = \text{DFT}^{-1}(\text{DFT}(f) \odot \text{DFT}(g))$

\begin{itemize}
  \item DFT naïve~: $O(n^2)$
  \item FFT (Fast Fourier Transform)~: $O(n \log n)$
\end{itemize}

%══════════════════════════════════════════════════════════════════════
\section{La NTT --- version modulaire}
%══════════════════════════════════════════════════════════════════════

\begin{defbox}{NTT (Number Theoretic Transform)}
La NTT est la DFT sur $\Zq$ au lieu de $\mathbb{C}$.
Elle remplace $\omega = e^{2\pi i / n}$ par une \textbf{racine primitive
$n$-ième de l'unité dans $\Zq$}, notée $\zeta$.

$$\hat{f}_k = \sum_{j=0}^{n-1} f_j \cdot \zeta^{jk} \pmod{q}$$

\textbf{Condition~:} $\zeta$ existe dans $\Zq$ si $n \mid (q-1)$.
\end{defbox}

\subsection{Les racines de l'unité dans $\Zq$}

\begin{importantbox}
Pour Dilithium~: $q = 8\,380\,417$ et $n = 256$.

On a $q - 1 = 8\,380\,416 = 2^{23} \times \text{impair}$.
Donc $2n = 512$ divise $q - 1$ --- la NTT de taille $n = 256$ existe !

La racine primitive $2n = 512$-ième de l'unité est~:
$$\zeta_{512} = 1753 \pmod{q}$$
On vérifie~: $1753^{512} \equiv 1 \pmod{q}$ et $1753^{256} \not\equiv 1$.
\end{importantbox}

\begin{lstlisting}[style=python]
q = 8_380_417
zeta = 1753  # racine primitive 512-ieme

# Verification
print(f"zeta^512 = {pow(zeta, 512, q)}")  # 1
print(f"zeta^256 = {pow(zeta, 256, q)}")  # != 1 (devrait etre -1)
print(f"zeta^256 + 1 = {(pow(zeta, 256, q) + 1) % q}")  # 0

# Generer les puissances de zeta (twiddle factors)
zetas = [pow(zeta, k, q) for k in range(512)]
print(f"Les 8 premiers zeta^k : {zetas[:8]}")
\end{lstlisting}

%══════════════════════════════════════════════════════════════════════
\section{NTT pour $R_q = \Zq[x]/(x^n+1)$}
%══════════════════════════════════════════════════════════════════════

\subsection{NTT "twisted" (négacyclique)}

La NTT standard s'applique au produit cyclique $\Zq[x]/(x^n-1)$.
Pour $\Rq = \Zq[x]/(x^n+1)$, on utilise une version "twisted" (négacyclique).

\begin{theorembox}{NTT négacyclique}
Pour $f \in \Rq$ de taille $n$, la NTT négacyclique est~:
$$\hat{f}_k = \sum_{j=0}^{n-1} f_j \cdot \psi^j \cdot \omega^{jk}, \quad k = 0, \ldots, n-1$$
où~:
\begin{itemize}
  \item $\psi = \zeta_{2n}$ est la racine primitive $2n$-ième (ici~: $\zeta_{512} = 1753$)
  \item $\omega = \psi^2 = \zeta_n$ est la racine $n$-ième
\end{itemize}
\end{theorembox}

\textbf{Conséquence~:} La multiplication dans $\Rq$ devient~:
$$\text{NTT}(f \cdot g \bmod x^n+1) = \text{NTT}(f) \odot \text{NTT}(g)$$

Ce qui donne l'algorithme de multiplication efficace~:
\begin{enumerate}
  \item $\hat{f} = \NTT(f)$ \hfill$O(n \log n)$
  \item $\hat{g} = \NTT(g)$ \hfill$O(n \log n)$
  \item $\hat{h} = \hat{f} \odot \hat{g}$ (pointwise) \hfill$O(n)$
  \item $h = \INTT(\hat{h})$ \hfill$O(n \log n)$
\end{enumerate}

\subsection{Dans le projet}

\begin{importantbox}
Dans le code Python du projet, les polynômes ont deux représentations~:

\begin{itemize}
  \item \textbf{Représentation normale}~: \texttt{poly.coeffs} (liste de 256 entiers)
  \item \textbf{Représentation NTT}~: \texttt{poly\_hat.coeffs} (256 entiers en domaine NTT)
\end{itemize}

Les conversions se font avec~:
\begin{itemize}
  \item \texttt{poly.to\_ntt()} : transforme en NTT (forward transform)
  \item \texttt{poly\_hat.from\_ntt()} : NTT inverse
\end{itemize}

Quand on voit \texttt{A\_hat @ s1.to\_ntt()}, c'est~:
\begin{enumerate}
  \item Convertir $s_1$ en NTT
  \item Multiplier terme-à-terme chaque ligne de $\hat{A}$ avec $\hat{s_1}$
  \item Sommer les produits (dans le domaine NTT)
\end{enumerate}
\end{importantbox}

%══════════════════════════════════════════════════════════════════════
\section{Algorithme Cooley-Tukey --- implémentation}
%══════════════════════════════════════════════════════════════════════

\subsection{NTT itérative (butterfly)}

L'algorithme NTT opère par ``butterfly''~: à chaque étape, on divise
le problème en deux moitiés et on combine les résultats.

\begin{lstlisting}[style=python]
def ntt_forward(f, q=8_380_417, n=256, zeta=1753):
    """
    NTT negacyclique iterative (Cooley-Tukey).
    Transforme in-place le tableau f.
    """
    f = f.copy()
    psi = zeta  # racine 2n-ieme

    k = 1
    length = n // 2
    while length >= 1:
        for start in range(0, n, 2 * length):
            w = pow(psi, k, q)  # twiddle factor
            k += 1
            for j in range(start, start + length):
                t = (w * f[j + length]) % q
                f[j + length] = (f[j] - t) % q
                f[j] = (f[j] + t) % q
        length //= 2

    return f

def ntt_inverse(f_hat, q=8_380_417, n=256, zeta=1753):
    """
    NTT inverse.
    """
    f = f_hat.copy()
    psi = zeta
    psi_inv = pow(psi, q - 2, q)  # inverse de psi
    n_inv = pow(n, q - 2, q)       # inverse de n

    k = n - 1
    length = 1
    while length <= n // 2:
        for start in range(0, n, 2 * length):
            w = pow(psi_inv, k, q)
            k -= 1
            for j in range(start, start + length):
                t = f[j]
                f[j] = (t + f[j + length]) % q
                f[j + length] = (w * (f[j + length] - t)) % q
        length *= 2

    # Normalisation par n^(-1)
    for i in range(n):
        f[i] = (n_inv * f[i]) % q

    return f

# Test de coherence
import random
q = 8_380_417
n = 256

poly = [random.randint(0, q - 1) for _ in range(n)]
poly_hat = ntt_forward(poly, q, n)
poly_back = ntt_inverse(poly_hat, q, n)

print(f"NTT coherente : {poly == poly_back}")
\end{lstlisting}

\subsection{Multiplication via NTT}

\begin{lstlisting}[style=python]
def poly_mul_ntt(f, g, q=8_380_417, n=256):
    """
    Multiplication dans R_q via NTT.
    3x NTT(n log n) + 1x produit pointe (n).
    """
    f_hat = ntt_forward(f, q, n)
    g_hat = ntt_forward(g, q, n)

    # Produit pointe
    fg_hat = [(f_hat[i] * g_hat[i]) % q for i in range(n)]

    # NTT inverse
    return ntt_inverse(fg_hat, q, n)

# Verification : comparer avec schoolbook
import random
q, n = 8_380_417, 8  # petit n pour tester
f = [random.randint(0, q-1) for _ in range(n)]
g = [random.randint(0, q-1) for _ in range(n)]

result_ntt = poly_mul_ntt(f, g, q, n)
result_schoolbook = poly_mul_schoolbook(f, g, q, n)

print(f"NTT == Schoolbook : {result_ntt == result_schoolbook}")
\end{lstlisting}

%══════════════════════════════════════════════════════════════════════
\section{Les twiddle factors dans le projet Solidity}
%══════════════════════════════════════════════════════════════════════

\begin{importantbox}
Dans \texttt{ZKNOX\_NTT.sol}, les twiddle factors (puissances de $\zeta$)
sont \textbf{précompilés} et stockés dans des contrats déployés séparément.

À l'exécution, \texttt{extcodecopy} lit les constantes depuis ces contrats.
Cela évite de les recalculer on-chain à chaque vérification.

C'est le composant le plus gas-critique : écrit entièrement en \textbf{Yul}
(assembleur EVM) pour maximiser l'efficacité.
\end{importantbox}

Voici comment les twiddle factors sont générés en Python pour être
exportés vers Solidity~:

\begin{lstlisting}[style=python]
def generate_twiddle_factors(q=8_380_417, n=256, zeta=1753):
    """
    Genere les twiddle factors pour la NTT de taille n.
    Ces constantes sont hardcodees dans ZKNOX_NTT.sol.
    """
    psi = zeta  # racine 2n-ieme
    twiddles = []

    k = 1
    length = n // 2
    while length >= 1:
        for _ in range(n // (2 * length)):
            twiddles.append(pow(psi, k, q))
            k += 1
        length //= 2

    return twiddles

twiddles = generate_twiddle_factors()
print(f"Nombre de twiddle factors : {len(twiddles)}")  # 255

# Exporter en Solidity (format uint256[])
solidity_array = ", ".join(str(t) for t in twiddles)
print(f"uint256[255] constant TWIDDLES = [{solidity_array[:100]}...]")
\end{lstlisting}

%══════════════════════════════════════════════════════════════════════
\section{Exercices}
%══════════════════════════════════════════════════════════════════════

\begin{exercice}{1 --- Vérification de la racine de l'unité}
Soit $q = 8\,380\,417$ et $\zeta = 1753$.
\begin{enumerate}[label=\alph*)]
  \item Vérifie que $\zeta^{512} \equiv 1 \pmod{q}$.
  \item Vérifie que $\zeta^{256} \equiv -1 \pmod{q}$.
  \item Calcule $\zeta^{128}$ --- est-ce une racine $4$-ième de l'unité ?
  \item Pourquoi $\zeta^{256} = -1$ est-il lié à $x^{256} + 1 = 0$ ?
\end{enumerate}
\end{exercice}

\begin{solution}
\begin{lstlisting}[style=python]
q, zeta = 8_380_417, 1753

# a)
print(pow(zeta, 512, q))   # 1  -> zeta est bien racine 512-ieme

# b)
print(pow(zeta, 256, q))   # q-1 = 8380416 = -1 mod q

# c)
val = pow(zeta, 128, q)
print(f"zeta^128 = {val}")
print(f"(zeta^128)^4 = {pow(val, 4, q)}")   # 1 -> c'est bien une racine 4-ieme

# d) L'anneau Z_q[x]/(x^256+1) impose x^256 = -1.
# La NTT "voit" zeta comme x : zeta^256 = -1 correspond exactement
# a la relation x^256 = -1 dans l'anneau.
\end{lstlisting}
\end{solution}

\begin{exercice}{2 --- NTT de taille 4}
Soit $q = 17$, $n = 4$, et $\zeta_8 = 2$ (racine primitive $8$-ième modulo $17$).

Vérifie~: $2^8 \equiv 1 \pmod{17}$ et $2^4 \equiv -1 \pmod{17}$.

\begin{enumerate}[label=\alph*)]
  \item Calcule $\NTT([1, 2, 3, 4])$ manuellement.
  \item Calcule $\INTT$ du résultat et vérifie que tu retrouves $[1,2,3,4]$.
\end{enumerate}
\end{exercice}

\begin{solution}
\begin{lstlisting}[style=python]
q, n, zeta = 17, 4, 2

# Verification
print(pow(zeta, 8, q))   # 1
print(pow(zeta, 4, q))   # 16 = -1 mod 17

# NTT forward
f = [1, 2, 3, 4]
f_hat = ntt_forward(f, q, n, zeta)
print(f"NTT([1,2,3,4]) = {f_hat}")

# NTT inverse
f_back = ntt_inverse(f_hat, q, n, zeta)
print(f"INTT(NTT(f)) = {f_back}")   # doit etre [1,2,3,4]
\end{lstlisting}
\end{solution}

\begin{exercice}{3 --- Comparaison de performances}
\begin{enumerate}[label=\alph*)]
  \item En Python, mesure le temps de \texttt{poly\_mul\_schoolbook} vs
        \texttt{poly\_mul\_ntt} pour $n = 256$.
        (Utilise \texttt{import time; t = time.perf\_counter()})
  \item Quel gain de vitesse obtiens-tu ?
  \item Pourquoi le gain théorique ($n^2$ vs $n\log n$) et le gain
        mesuré peuvent différer ?
\end{enumerate}
\end{exercice}

\begin{solution}
\begin{lstlisting}[style=python]
import time, random

q, n = 8_380_417, 256
f = [random.randint(0, q-1) for _ in range(n)]
g = [random.randint(0, q-1) for _ in range(n)]

# Schoolbook
t0 = time.perf_counter()
for _ in range(10):
    poly_mul_schoolbook(f, g, q, n)
t_school = (time.perf_counter() - t0) / 10

# NTT
t0 = time.perf_counter()
for _ in range(10):
    poly_mul_ntt(f, g, q, n)
t_ntt = (time.perf_counter() - t0) / 10

print(f"Schoolbook : {t_school*1000:.2f} ms")
print(f"NTT        : {t_ntt*1000:.2f} ms")
print(f"Gain       : {t_school/t_ntt:.1f}x")

# Note : le gain mesure < n/log(n) car :
# - Overhead Python (interpretation)
# - pow() sur grands entiers est lent en Python
# Le vrai gain apparait en C/Rust/Yul
\end{lstlisting}
\end{solution}

\begin{exercice}{4 --- Le domaine NTT dans le code}
Dans \texttt{Falcon\_dilithium.py}, la ligne~:
\begin{center}
\texttt{t = (A\_hat @ s1.to\_ntt()).from\_ntt() + s2}
\end{center}
\begin{enumerate}[label=\alph*)]
  \item Pourquoi convertit-on $s_1$ avec \texttt{.to\_ntt()} avant de multiplier par $A$ ?
  \item $\hat{A}$ est-il déjà en représentation NTT ? Comment est-il généré ?
  \item Pourquoi appelle-t-on \texttt{.from\_ntt()} avant d'ajouter $s_2$ ?
  \item Quel serait le problème si on oubliait le \texttt{.from\_ntt()} ?
\end{enumerate}
\end{exercice}

\end{document}
